\documentclass[12pt]{article}
\usepackage{graphicx} % Required for inserting images
\usepackage{fancyhdr}
\usepackage{enumitem}
\usepackage[margin=1in]{geometry}
\usepackage{amsmath}
\usepackage{amsfonts}
\usepackage{amsthm}
\usepackage{amssymb}
\usepackage{polynom}
\usepackage{mathrsfs}
\usepackage{xcolor}
\usepackage{parskip}
\usepackage{hyperref}
\usepackage{cancel}
\usepackage{slashed}
\usepackage{breqn}
\usepackage{enumitem}

\pagestyle{fancy}

\setlength{\headheight}{14.49998pt}
\setlength\parindent{0pt} %NO INDENT
\hypersetup{
	colorlinks=true, %set true if you want colored links
	linktoc=all,     %set to all if you want both sections and subsections linked
	linkcolor=blue,  %choose some color if you want links to stand out
}

\DeclareMathOperator*{\res}{Res}
\DeclareMathOperator{\Div}{div}
\DeclareMathOperator{\curl}{curl}
\DeclareMathOperator{\Int}{int}
\DeclareMathOperator{\diam}{diam}
\DeclareMathOperator{\range}{range}
\DeclareMathOperator{\grad}{\nabla}
\DeclareMathOperator{\Grad}{grad}
\DeclareMathOperator{\supp}{supp}
\DeclareMathOperator{\dist}{dist}
\DeclareMathOperator{\Span}{span}

\newtheorem{proposition}{Proposition}[section]
\newtheorem{theorem}[proposition]{Theorem}
\newtheorem{corollary}[proposition]{Corollary}
\newtheorem{lemma}[proposition]{Lemma}
\newtheorem{definition}[proposition]{Definition}
\newcommand{\RR}{\mathbb{R}}
\newcommand{\nablaa}{\cancel{\nabla}}
\newcommand{\lec}{\lesssim}
\newenvironment{solution}{\begin{proof}[Solution]}{\end{proof}}
\newcommand{\oo}{\"{o}}

\title{Homework 3}
\author{Zucca}

\fancypagestyle{firstpage}{
	\fancyhf{}
	\fancyhead[L]{PDEs Homework 5}
	\fancyhead[R]{Steve Zucca}
	\fancyfoot[C]{\thepage}
	\renewcommand{\headrulewidth}{0.4pt}
}
\fancyhead[R]{Zucca}
\begin{document}
	\thispagestyle{firstpage}
	\textbf{Problem \# 1}. (Exercise 1.4.1a) Let $H(t,x)$ be the heat kernel on $(0,\infty)\times \mathbb{R}^n$. Define its extension to $\mathbb{R}\times \mathbb{R}^n$ by 
	\[
		H_e(t,x)=
		\begin{cases}
			H(t,x) & : t> 0 \\
			0 & : t\leq 0
		\end{cases}
	\]
	\begin{enumerate}[label=(\alph*)]
		\item Show that $H_e \in L^1_{\text{loc}}(\mathbb{R}\times \mathbb{R}^n)$, and thus $H_e\in \mathcal{D}'(\mathbb{R}\times \mathbb{R}^n)$.

	\begin{solution} Let $K\subset \mathbb{R}\times \mathbb{R}^n$ be compact. So, $K\subset [-T,T]\times\overline{B(0,R)}$, where $0<T,R<\infty$. Therefore,
		\begin{equation}
			\begin{aligned}
				\int_K |H_e|\,dm &= \int_{-T}^T \int_{\overline{B(0,R)}} |H_e|\,dx\,dt = \lim_{\epsilon\to 0^+} \int_\epsilon^T \int_{\overline{B(0,R)}}|H(t,x)|\,dx\,dt \\ 
				&= \lim_{\epsilon\to 0^+}\int_\epsilon^T \int_{\overline{B(0,R)}} (4\pi t)^{-\frac{n}{2}}e^{-\frac{|x|^2}{4t}}\,dx\,dt \\
				&\leq \lim_{\epsilon\to 0^+} \int_\epsilon^T (4\pi t)^{-\frac{n}{2}}\int_{\mathbb{R}^n}e^{-\frac{|x|^2}{4t}}\,dx\,dt = \lim_{\epsilon\to 0^+}\int_\epsilon^T (4\pi t)^{-\frac{n}{2}}(4\pi t)^{\frac{n}{2}}\,dt \\
				&= \lim_{\epsilon\to 0^+} T-\epsilon = T \\
				&< \infty
			\end{aligned}
		\end{equation}
		So, $H_e\in L^1_{\text{loc}}(\mathbb{R}\times \mathbb{R}^n)$.
	\end{solution}
\item Show that $ H_e $ is a fundamental solution for the heat operator, that is,
\[
	(\partial_t -\Delta)H_e = \delta\in \mathcal{D}'(\mathbb{R}\times \mathbb{R}^n)
\]
\begin{solution}
	For $ \eta > 0 $, define
	\[
		H_\eta(x,t) = \begin{cases}
			H_e & : t> \eta \\
			0 & : t\leq\eta
		\end{cases}
	\]
	\begin{enumerate}[label=(\roman*)]
		\item $H_\eta\to H_e$ in $\mathcal{D}'(\mathbb{R}\times \mathbb{R}^n)$, as $\eta\to 0$.
		
		Let $\phi\in \mathcal{D}(\mathbb{R}\times \mathbb{R}^n)$, and let $\epsilon > 0$. Choose $\eta>0$ so that $\eta < \frac{\epsilon}{\left|\left|\phi \right|\right|_{L^\infty}}$. Assume that $\supp \phi\subset [-T,T]\times\overline{B(0,R)}$, where $0<T,R<\infty$. Then,
		\begin{equation}
			\begin{aligned}
				\left| \langle H_\eta,\phi\rangle - \langle H_e,\phi\rangle \right|&= \left| \int (H_\eta-H_e)\phi\,dm \right| = \left| \int_{0}^\eta \int_{\overline{B(0,R)}} -H_e\phi\,dx\,dt \right| \\
				&\leq \left|\left|\phi \right|\right|_{L^\infty} \int_{0}^\eta \int_{\mathbb{R}^n} H_e\,dx\,dt = \left|\left|\phi \right|\right|_{L^\infty} \lim_{\delta\to 0^+} \int_{\delta}^\eta \int_{\mathbb{R}^n} H\,dx\,dt \\
				&= \left|\left|\phi \right|\right|_{L^\infty} \eta \\
				& < \epsilon
			\end{aligned}
		\end{equation}
		So, $H_\eta\to H_e$ in as $\eta\to 0$.
		\newpage 
		\item $(\partial_t-\Delta)H_\eta \to \delta$ in $\mathcal{D}'(\mathbb{R}\times \mathbb{R}^n)$ as $\eta\to 0$.
		
		For $\phi\in \mathcal{D}(\mathbb{R}\times \mathbb{R}^n)$, we have
		\begin{equation}
			\begin{aligned}
				\langle (\partial_t-\Delta)H_\eta,\phi\rangle &= \langle H_\eta,(-\partial_t-\Delta)\phi\rangle = \int_{\eta}^\infty \int_{\mathbb{R}^n} H_e(-\partial_t-\Delta)\phi\,dx\,dt \\
				&= \int_{\mathbb{R}^n} H(\eta,x)\phi(\eta,x)\,dx + \cancel{\int_{\eta}^\infty \int_{\mathbb{R}^n} (\partial_t-\Delta)H\phi\,dx\,dt}
			\end{aligned}
		\end{equation}
		Since $H$ is an approximate identity, we have that 
		\begin{equation}\label{Delta}
			\lim_{\eta\to 0} \langle (\partial_t-\Delta)H_\eta,\phi\rangle = \phi(0) = \langle \delta,\phi\rangle
		\end{equation}
		Since this holds for all such $\phi$, we have that $\lim_{\eta\to 0}(\partial_t-\Delta)H_\eta = \delta$ in $\mathcal{D}'(\mathbb{R}\times \mathbb{R}^n).$
		
		Clearly, $(\partial_t-\Delta)$ is a differential operator with smooth coefficients. Therefore, $(\partial_t-\Delta):\mathcal{D}'(\mathbb{R}\times \mathbb{R}^n)\to \mathcal{D}'(\mathbb{R}\times \mathbb{R}^n)$ is continuous, and we have
		\begin{equation}
			\delta = \lim_{\eta\to 0}(\partial_t-\Delta )H_\eta = (\partial_t-\Delta)(\lim_{\eta\to 0}H_\eta) = (\partial_t-\Delta)H_e
		\end{equation}
		\item Where does your argument use the fact that the topology of $\mathcal{D}'(\mathbb{R}\times \mathbb{R}^n)$ is Hausdorff?
		
		We used Hausdorff(ness) in (5). We have that $\lim_{\eta\to 0}(\partial_t-\Delta)H_\eta = \delta$, and we have that $\lim_{\eta\to 0}(\partial_t-\Delta)H_\eta = (\partial_t-\Delta)H_e$. We use the fact that limits are unique to conclude that $\delta=(\partial_t-\Delta)H_e$.
	\end{enumerate}
\end{solution}
		\end{enumerate}
	\textbf{Problem \# 4}. (\S 1.7 Problem 2) Consider the Schrodinger initial value problem
	\[
		\begin{cases}
			u_t-i\Delta u = 0 \\
			u(0,\cdot) = g\in \mathcal{S}(\mathbb{R}^n)
		\end{cases}
	\]
	\begin{enumerate}[label=(\alph*)]
		\item Starting with the heat equation, explain how a complex change of variable $t\mapsto it$ can be used to quickly derive the solution formula
		\[
			u(t,x)=(4\pi i t)^{-\frac{n}{2}}\int e^{-\frac{|x-y|^2}{4it}}g(y)\,dy
		\]
		\begin{solution}
			Suppose $w$ solves 
			\begin{equation}
				\begin{cases}
					w_t-\Delta w = 0 \\
					w(0,\cdot)=g\in \mathcal{S}(\mathbb{R}^n)
				\end{cases}
			\end{equation}
			We know that the \textit{unique} solution to this problem is 
			\begin{equation}
				w(t,x)=(4\pi t)^{-\frac{n}{2}}\int e^{-\frac{|x-y|^2}{4t}}g(y)\,dy
			\end{equation}
			Now, let $s= -it$, and let $\tilde{w}=w(s,x)$. By the chain rule, we know $\partial_t\mapsto -i\partial_s$, so we have
			\begin{equation}
				\begin{cases}
					-i\tilde{w}_s-\Delta\tilde{w}=0\implies \tilde{w}_s-i\Delta\tilde{w}=0 \\
					\tilde{w}(0,\cdot)=g
				\end{cases}
			\end{equation}
			In other words, $\tilde{w}$ solves the Schrodinger initial value problem. Since all we did was change variables, we have that
			\begin{equation}
				\tilde{w}(t,x)=w(it,x) = (4\pi it)^{-\frac{n}{2}}\int e^{-\frac{|x-y|^2}{4it}}g(y)\,dy
			\end{equation}
			as desired. This is a \textit{derivation}. To actually show that this candidate solves the IVP, we have to go through the usual steps to justify differentiation under the integral sign.
		\end{solution}
		\item Prove the dispersive estimate
		\[
			\left| u(t,x) \right|\leq C t^{-\frac{n}{2}}\left|\left|g \right|\right|_{L^1} 
		\]
		\begin{solution}
			This is trivial.
			\[
				\left| u(t,x) \right| \leq \left|(4\pi i t)^{-\frac{n}{2}}\right| \int \left|e^{-\frac{|x-y|^2}{4it}}g(y)\right| \,dy = |(4\pi i t)|^{-\frac{n}{2}}\int 1\cdot |g(y)|\,dy \leq C t^{-\frac{n}{2}}\left|\left|g \right|\right|_{L^1}
			\]
		\end{solution}
		\item Define $E(t)=\int |u|^2\,dx$. Use an energy method to show that $E$ is a constant function of $t$.
		\begin{solution}
			\begin{align*}
				E'(t) &= \int u_t \overline{u} + u \overline{u}_t\,dx = 2\Re \int u_t\overline{u}\,dx = 2\Re \int i\Delta u \overline{u}\,dx \\
				&= 2\Re \int (-i) |\nabla u|^2\,dx = 0
			\end{align*}
			since the last quantity is pure imaginary. To get to the second line, we used Gauss-Green and the fact that the solution $u\in S(\mathbb{R}_+\times \mathbb{R}^n)$, so there are no boundary terms.
			
			Therefore, $E$ is a constant function of $t$.
		\end{solution}
		\item Re-derive the result of (c) using the Fourier Transform.
		\begin{solution}
			By Plancherel's Theorem\footnote{This was proven in Harmonic Analysis. The key property is that $\int \hat{f}g\,d'\xi=\int f\hat{g}\,d'x$.}
			\begin{equation}
				\begin{aligned}
					\left|\left|u \right|\right|_{L^2_x}^2 &= \left|\left|\hat{u} \right|\right|_{L^2_\xi}^2\\ &= \left|\left|\int e^{-ix\cdot\xi}u(t,x)\,d'x \right|\right|_{L^2_\xi}^2\\ &= \left|\left|(4\pi i t)^{-\frac{n}{2}}\int e^{-ix\cdot\xi}\int e^{-\frac{|x-y|^2}{4it}}g(y)\,dy\,d'x \right|\right|_{L^2_\xi}^2 \\
					&\stackrel{A}{=} \left|\left|(4\pi i t)^{-\frac{n}{2}}\int g(y)\int e^{-ix\cdot\xi}e^{-\frac{|x-y|^2}{4it}}\,d'x\,dy \right|\right|_{L^2_\xi}^2 \\
					&= \left|\left|(4\pi i t)^{-\frac{n}{2}}\int g(y) e ^{-iy\cdot\xi}\int e^{-ix\cdot\xi}e^{-\frac{|x|^2}{4it}}\,d'x\,dy \right|\right|_{L^2_\xi}^2&\qquad (\text{F.T. of Translation}) \\
					&\stackrel{B}{=} \left|\left|(4\pi i t)^{-\frac{n}{2}}\int g(y) e ^{-iy\cdot\xi} (2it)^{\frac{n}{2}}\int e^{-ix\cdot(\sqrt{2it}\xi)}e^{-\frac{|x|^2}{2}}\,d'x\,dy \right|\right|_{L^2_\xi}^2 \\
					&= \left|\left|(4\pi i t)^{-\frac{n}{2}}\int g(y) e ^{-iy\cdot\xi} (2it)^{\frac{n}{2}}\int e^{-ix\cdot(\sqrt{2it}\xi)}e^{-\frac{|x|^2}{2}}\,d'x\,dy \right|\right|_{L^2_\xi}^2 \\
					&= \left|\left|(2\pi)^{-\frac{n}{2}}e^{-it|\xi|^2}\int g(y) e ^{-iy\cdot\xi} \,dy \right|\right|_{L^2_\xi}^2 \\
					&= \left|\left|\int g(y)e^{-iy\cdot\xi}\,d'y \right|\right|_{L^2_\xi}^2 &\qquad \text{since } \left| e^{-it|\xi|^2} \right|=1  \\
					&= \left|\left|\hat{g}(\xi) \right|\right|_{L^2_\xi}^2 \\
					&= \left|\left|g \right|\right|_{L^2_x}^2
				\end{aligned}
			\end{equation}
			In equality $(A)$, we used Fubini's Theorem which is valid since $g\in \mathbb{S}(\mathbb{R}^n)$ and the Gaussian is rapidly decaying.
			
			In equality $(B)$, we are using Cauchy's Residue Theorem. We only have the Fourier dilation theorem for real scalars, so we have to justify deforming the contour. I'd rather not get into this.
			
			Clearly, this result is independent of $t$.
		\end{solution}
	\end{enumerate}
	
\end{document}